% !TEX program = lualatex
\ifdefined\pg\else\IfFileExists{tex/avulso.tex}{% Compilação AVULSA de um fragmento (volN/.../<stem>.tex) — p. ex. pelo botão do editor,
% de dentro da pasta do nó.  O fragmento começa com
%
%   \ifdefined\pg\else\IfFileExists{tex/avulso.tex}{% Compilação AVULSA de um fragmento (volN/.../<stem>.tex) — p. ex. pelo botão do editor,
% de dentro da pasta do nó.  O fragmento começa com
%
%   \ifdefined\pg\else\IfFileExists{tex/avulso.tex}{% Compilação AVULSA de um fragmento (volN/.../<stem>.tex) — p. ex. pelo botão do editor,
% de dentro da pasta do nó.  O fragmento começa com
%
%   \ifdefined\pg\else\IfFileExists{tex/avulso.tex}{\input{tex/avulso}}{\input{../../tex/avulso}}\fi
%
% e termina com  \csname fimavulso\endcsname .  Dentro do livro ou do _preview, \pg já existe
% e nada disto roda; \fimavulso fica indefinido e o \csname vira \relax.
%
% Motor: o projeto é lualatex (fontspec/polyglossia).  Se o editor chamar pdflatex, cai
% num preâmbulo equivalente com inputenc/babel — serve para conferir, não para o PDF final.
\documentclass[11pt,a5paper]{book}
\ifdefined\directlua
  \usepackage{fontspec}
  \usepackage{polyglossia}
  \setmainlanguage{portuguese}
\else\ifdefined\XeTeXrevision
  \usepackage{fontspec}
  \usepackage{polyglossia}
  \setmainlanguage{portuguese}
\else
  \usepackage[utf8]{inputenc}
  \usepackage[T1]{fontenc}
  \usepackage{lmodern}
  \usepackage[portuguese]{babel}
  \DeclareUnicodeCharacter{2192}{\ensuremath{\rightarrow}}
\fi\fi
% raiz do repo: compilado da raiz ou de volN/<parte>/ (dois níveis abaixo)
\IfFileExists{tex/preambulo.tex}{\input{tex/preambulo}}{\input{../../tex/preambulo}}
\def\fimavulso{\end{document}}
\begin{document}
}{% Compilação AVULSA de um fragmento (volN/.../<stem>.tex) — p. ex. pelo botão do editor,
% de dentro da pasta do nó.  O fragmento começa com
%
%   \ifdefined\pg\else\IfFileExists{tex/avulso.tex}{\input{tex/avulso}}{\input{../../tex/avulso}}\fi
%
% e termina com  \csname fimavulso\endcsname .  Dentro do livro ou do _preview, \pg já existe
% e nada disto roda; \fimavulso fica indefinido e o \csname vira \relax.
%
% Motor: o projeto é lualatex (fontspec/polyglossia).  Se o editor chamar pdflatex, cai
% num preâmbulo equivalente com inputenc/babel — serve para conferir, não para o PDF final.
\documentclass[11pt,a5paper]{book}
\ifdefined\directlua
  \usepackage{fontspec}
  \usepackage{polyglossia}
  \setmainlanguage{portuguese}
\else\ifdefined\XeTeXrevision
  \usepackage{fontspec}
  \usepackage{polyglossia}
  \setmainlanguage{portuguese}
\else
  \usepackage[utf8]{inputenc}
  \usepackage[T1]{fontenc}
  \usepackage{lmodern}
  \usepackage[portuguese]{babel}
  \DeclareUnicodeCharacter{2192}{\ensuremath{\rightarrow}}
\fi\fi
% raiz do repo: compilado da raiz ou de volN/<parte>/ (dois níveis abaixo)
\IfFileExists{tex/preambulo.tex}{\input{tex/preambulo}}{\input{../../tex/preambulo}}
\def\fimavulso{\end{document}}
\begin{document}
}\fi
%
% e termina com  \csname fimavulso\endcsname .  Dentro do livro ou do _preview, \pg já existe
% e nada disto roda; \fimavulso fica indefinido e o \csname vira \relax.
%
% Motor: o projeto é lualatex (fontspec/polyglossia).  Se o editor chamar pdflatex, cai
% num preâmbulo equivalente com inputenc/babel — serve para conferir, não para o PDF final.
\documentclass[11pt,a5paper]{book}
\ifdefined\directlua
  \usepackage{fontspec}
  \usepackage{polyglossia}
  \setmainlanguage{portuguese}
\else\ifdefined\XeTeXrevision
  \usepackage{fontspec}
  \usepackage{polyglossia}
  \setmainlanguage{portuguese}
\else
  \usepackage[utf8]{inputenc}
  \usepackage[T1]{fontenc}
  \usepackage{lmodern}
  \usepackage[portuguese]{babel}
  \DeclareUnicodeCharacter{2192}{\ensuremath{\rightarrow}}
\fi\fi
% raiz do repo: compilado da raiz ou de volN/<parte>/ (dois níveis abaixo)
\IfFileExists{tex/preambulo.tex}{% Preâmbulo comum ao livro e à pré-visualização de fragmentos (tex/_preview.tex).

% Mancha em A5.  Sem isto vale o padrão da classe book (~105 mm de largura e margem
% externa de ~25 mm em frente e verso), largo demais para o formato.  Aqui: ~118 mm de
% largura, margem interna um pouco maior que a externa (encadernação).
\usepackage[a5paper,inner=16mm,outer=14mm,top=14mm,bottom=16mm,
            includehead,headsep=5mm,footskip=8mm]{geometry}

% Títulos de secção sem justificar (evita linhas estouradas em \Large bold, p. ex.
% «Demonstrabilidade») e um pouco de folga para enunciados de tese em negrito.
\usepackage{sectsty}
\allsectionsfont{\raggedright}
\setlength{\emergencystretch}{1.5em}

% Marca da página impressa do original: \pg{286} imprime [286] no fluxo do texto.
\newcommand{\pg}[1]{\,{\footnotesize\textbf{[#1]}}\,}

% Virada de página do original DENTRO de uma nota (\pgn{422}); discreta, para não quebrar a nota.
\newcommand{\pgn}[1]{{\scriptsize\textbf{[#1]}}}

% Asterisco do original (divisão de assunto, sem título): \asterisco centrado em linha própria,
% também válido no meio do parágrafo ou dentro de nota (\footnote{... \asterisco ...}).
\newcommand{\asterisco}{\par\medskip\centerline{\footnotesize$\ast$}\medskip\noindent}

% Gravura recortada do fac-símile (<figure> do TEI; tools/tei2latex.py): fora de float, no ponto
% exato do original.  \gravura[detalhes]{arquivo.png}{fração da largura do texto}{legenda}:
% legenda principal centrada ("FIGURA VIII. — Côrte…"); os detalhes opcionais (<p> do <figure>:
% "A, hemispherio…; — B, …") vão abaixo, em corpo menor e justificados.  Ambos podem ser vazios.
\usepackage{graphicx}
\newcommand{\gravura}[4][]{\par\medskip\begingroup\centering
  \includegraphics[width=#3\linewidth,height=0.75\textheight,keepaspectratio]{#2}\par
  \if\relax\detokenize{#4}\relax\else\smallskip{\small #4\par}\fi
  \if\relax\detokenize{#1}\relax\else\smallskip{\footnotesize\leftskip=1em\rightskip=1em
    \parfillskip=0pt plus 1fil\noindent #1\par}\fi\endgroup\medskip\noindent}

% Citações latinas de S. Tomás (o OCR erra muito aqui: sempre conferir na imagem).
\newenvironment{latim}{\begin{quote}\itshape}{\end{quote}}

% --- Estrutura do autor (promovido de "Artigo III.tex") -----------------------
% enunciado de tese numerada que corre com o texto (segue-se travessão)
\newcommand{\tese}[2]{\bigskip\noindent\textbf{#1.\quad #2}\;\textemdash\;}
% enunciado de tese numerada em linha própria (seguem-se as subdivisões a), b)...)
\newcommand{\teseB}[2]{\bigskip\noindent\textbf{#1.\quad #2}\par}
% subdivisão a), b), c)... com rubrica em itálico
\newcommand{\sub}[2]{\medskip\noindent#1)\quad\emph{#2},\;\textemdash\;}
% subdivisão a), b), c)... sem rubrica
\newcommand{\subp}[1]{\medskip\noindent#1)\quad}

% --- Notas: no pé (padrão) ou correndo no corpo do texto ---------------------
% Passar  \def\notascorridas{}  no -usepretex faz cada nota correr no corpo, em
% corpo menor, num bloco recuado que interrompe a leitura até acabar.  Ex.:
%   latexmk -lualatex -usepretex='\def\fragmento{vol2/.../01_cap_y.tex}\def\notascorridas{}' tex/_preview.tex
\newif\ifnotascorridas
\ifdefined\notascorridas \notascorridastrue \fi
\ifnotascorridas
  \renewcommand\footnote[1]{%
    \stepcounter{footnote}\textsuperscript{\thefootnote}% marca no ponto do texto
    \par
    {\footnotesize
     \leftskip=2em \rightskip=2em \parindent=1em
     \addvspace{\medskipamount}%
     \noindent\textit{Nota \thefootnote.}\enspace #1\par
     \addvspace{\medskipamount}}%
  }
\fi

% --- Sem notas: \def\semnotas{} no -usepretex suprime todas as notas ---------
% (nem a chamada no texto nem o conteúdo).  Usado por  ./compila.sh --sem-notas .
\ifdefined\semnotas
  \renewcommand\footnote[2][]{}%
  \renewcommand\footnotemark[1][]{}%
  \renewcommand\footnotetext[2][]{}%
\fi

% peças pré-textuais (nó pretextual; p@type em web/obra/esquema.py)
\newcommand{\preTituloObra}[1]{\par\vspace{1em}{\centering\Large\bfseries\scshape #1\par}\vspace{1em}}
\newcommand{\preCentro}[1]{\par{\centering #1\par}}
\newcommand{\preEnderecamento}[1]{\par\vspace{1em}{\centering\small\scshape #1\par}}
\newcommand{\preSaudacao}[1]{\par\medskip\noindent #1\par}
\newcommand{\preDataLocal}[1]{\par\medskip\noindent\emph{#1}\par}
\newcommand{\preFecho}[1]{\par\smallskip{\raggedleft\scshape #1\par}}
}{% Preâmbulo comum ao livro e à pré-visualização de fragmentos (tex/_preview.tex).

% Mancha em A5.  Sem isto vale o padrão da classe book (~105 mm de largura e margem
% externa de ~25 mm em frente e verso), largo demais para o formato.  Aqui: ~118 mm de
% largura, margem interna um pouco maior que a externa (encadernação).
\usepackage[a5paper,inner=16mm,outer=14mm,top=14mm,bottom=16mm,
            includehead,headsep=5mm,footskip=8mm]{geometry}

% Títulos de secção sem justificar (evita linhas estouradas em \Large bold, p. ex.
% «Demonstrabilidade») e um pouco de folga para enunciados de tese em negrito.
\usepackage{sectsty}
\allsectionsfont{\raggedright}
\setlength{\emergencystretch}{1.5em}

% Marca da página impressa do original: \pg{286} imprime [286] no fluxo do texto.
\newcommand{\pg}[1]{\,{\footnotesize\textbf{[#1]}}\,}

% Virada de página do original DENTRO de uma nota (\pgn{422}); discreta, para não quebrar a nota.
\newcommand{\pgn}[1]{{\scriptsize\textbf{[#1]}}}

% Asterisco do original (divisão de assunto, sem título): \asterisco centrado em linha própria,
% também válido no meio do parágrafo ou dentro de nota (\footnote{... \asterisco ...}).
\newcommand{\asterisco}{\par\medskip\centerline{\footnotesize$\ast$}\medskip\noindent}

% Gravura recortada do fac-símile (<figure> do TEI; tools/tei2latex.py): fora de float, no ponto
% exato do original.  \gravura[detalhes]{arquivo.png}{fração da largura do texto}{legenda}:
% legenda principal centrada ("FIGURA VIII. — Côrte…"); os detalhes opcionais (<p> do <figure>:
% "A, hemispherio…; — B, …") vão abaixo, em corpo menor e justificados.  Ambos podem ser vazios.
\usepackage{graphicx}
\newcommand{\gravura}[4][]{\par\medskip\begingroup\centering
  \includegraphics[width=#3\linewidth,height=0.75\textheight,keepaspectratio]{#2}\par
  \if\relax\detokenize{#4}\relax\else\smallskip{\small #4\par}\fi
  \if\relax\detokenize{#1}\relax\else\smallskip{\footnotesize\leftskip=1em\rightskip=1em
    \parfillskip=0pt plus 1fil\noindent #1\par}\fi\endgroup\medskip\noindent}

% Citações latinas de S. Tomás (o OCR erra muito aqui: sempre conferir na imagem).
\newenvironment{latim}{\begin{quote}\itshape}{\end{quote}}

% --- Estrutura do autor (promovido de "Artigo III.tex") -----------------------
% enunciado de tese numerada que corre com o texto (segue-se travessão)
\newcommand{\tese}[2]{\bigskip\noindent\textbf{#1.\quad #2}\;\textemdash\;}
% enunciado de tese numerada em linha própria (seguem-se as subdivisões a), b)...)
\newcommand{\teseB}[2]{\bigskip\noindent\textbf{#1.\quad #2}\par}
% subdivisão a), b), c)... com rubrica em itálico
\newcommand{\sub}[2]{\medskip\noindent#1)\quad\emph{#2},\;\textemdash\;}
% subdivisão a), b), c)... sem rubrica
\newcommand{\subp}[1]{\medskip\noindent#1)\quad}

% --- Notas: no pé (padrão) ou correndo no corpo do texto ---------------------
% Passar  \def\notascorridas{}  no -usepretex faz cada nota correr no corpo, em
% corpo menor, num bloco recuado que interrompe a leitura até acabar.  Ex.:
%   latexmk -lualatex -usepretex='\def\fragmento{vol2/.../01_cap_y.tex}\def\notascorridas{}' tex/_preview.tex
\newif\ifnotascorridas
\ifdefined\notascorridas \notascorridastrue \fi
\ifnotascorridas
  \renewcommand\footnote[1]{%
    \stepcounter{footnote}\textsuperscript{\thefootnote}% marca no ponto do texto
    \par
    {\footnotesize
     \leftskip=2em \rightskip=2em \parindent=1em
     \addvspace{\medskipamount}%
     \noindent\textit{Nota \thefootnote.}\enspace #1\par
     \addvspace{\medskipamount}}%
  }
\fi

% --- Sem notas: \def\semnotas{} no -usepretex suprime todas as notas ---------
% (nem a chamada no texto nem o conteúdo).  Usado por  ./compila.sh --sem-notas .
\ifdefined\semnotas
  \renewcommand\footnote[2][]{}%
  \renewcommand\footnotemark[1][]{}%
  \renewcommand\footnotetext[2][]{}%
\fi

% peças pré-textuais (nó pretextual; p@type em web/obra/esquema.py)
\newcommand{\preTituloObra}[1]{\par\vspace{1em}{\centering\Large\bfseries\scshape #1\par}\vspace{1em}}
\newcommand{\preCentro}[1]{\par{\centering #1\par}}
\newcommand{\preEnderecamento}[1]{\par\vspace{1em}{\centering\small\scshape #1\par}}
\newcommand{\preSaudacao}[1]{\par\medskip\noindent #1\par}
\newcommand{\preDataLocal}[1]{\par\medskip\noindent\emph{#1}\par}
\newcommand{\preFecho}[1]{\par\smallskip{\raggedleft\scshape #1\par}}
}
\def\fimavulso{\end{document}}
\begin{document}
}{% Compilação AVULSA de um fragmento (volN/.../<stem>.tex) — p. ex. pelo botão do editor,
% de dentro da pasta do nó.  O fragmento começa com
%
%   \ifdefined\pg\else\IfFileExists{tex/avulso.tex}{% Compilação AVULSA de um fragmento (volN/.../<stem>.tex) — p. ex. pelo botão do editor,
% de dentro da pasta do nó.  O fragmento começa com
%
%   \ifdefined\pg\else\IfFileExists{tex/avulso.tex}{\input{tex/avulso}}{\input{../../tex/avulso}}\fi
%
% e termina com  \csname fimavulso\endcsname .  Dentro do livro ou do _preview, \pg já existe
% e nada disto roda; \fimavulso fica indefinido e o \csname vira \relax.
%
% Motor: o projeto é lualatex (fontspec/polyglossia).  Se o editor chamar pdflatex, cai
% num preâmbulo equivalente com inputenc/babel — serve para conferir, não para o PDF final.
\documentclass[11pt,a5paper]{book}
\ifdefined\directlua
  \usepackage{fontspec}
  \usepackage{polyglossia}
  \setmainlanguage{portuguese}
\else\ifdefined\XeTeXrevision
  \usepackage{fontspec}
  \usepackage{polyglossia}
  \setmainlanguage{portuguese}
\else
  \usepackage[utf8]{inputenc}
  \usepackage[T1]{fontenc}
  \usepackage{lmodern}
  \usepackage[portuguese]{babel}
  \DeclareUnicodeCharacter{2192}{\ensuremath{\rightarrow}}
\fi\fi
% raiz do repo: compilado da raiz ou de volN/<parte>/ (dois níveis abaixo)
\IfFileExists{tex/preambulo.tex}{\input{tex/preambulo}}{\input{../../tex/preambulo}}
\def\fimavulso{\end{document}}
\begin{document}
}{% Compilação AVULSA de um fragmento (volN/.../<stem>.tex) — p. ex. pelo botão do editor,
% de dentro da pasta do nó.  O fragmento começa com
%
%   \ifdefined\pg\else\IfFileExists{tex/avulso.tex}{\input{tex/avulso}}{\input{../../tex/avulso}}\fi
%
% e termina com  \csname fimavulso\endcsname .  Dentro do livro ou do _preview, \pg já existe
% e nada disto roda; \fimavulso fica indefinido e o \csname vira \relax.
%
% Motor: o projeto é lualatex (fontspec/polyglossia).  Se o editor chamar pdflatex, cai
% num preâmbulo equivalente com inputenc/babel — serve para conferir, não para o PDF final.
\documentclass[11pt,a5paper]{book}
\ifdefined\directlua
  \usepackage{fontspec}
  \usepackage{polyglossia}
  \setmainlanguage{portuguese}
\else\ifdefined\XeTeXrevision
  \usepackage{fontspec}
  \usepackage{polyglossia}
  \setmainlanguage{portuguese}
\else
  \usepackage[utf8]{inputenc}
  \usepackage[T1]{fontenc}
  \usepackage{lmodern}
  \usepackage[portuguese]{babel}
  \DeclareUnicodeCharacter{2192}{\ensuremath{\rightarrow}}
\fi\fi
% raiz do repo: compilado da raiz ou de volN/<parte>/ (dois níveis abaixo)
\IfFileExists{tex/preambulo.tex}{\input{tex/preambulo}}{\input{../../tex/preambulo}}
\def\fimavulso{\end{document}}
\begin{document}
}\fi
%
% e termina com  \csname fimavulso\endcsname .  Dentro do livro ou do _preview, \pg já existe
% e nada disto roda; \fimavulso fica indefinido e o \csname vira \relax.
%
% Motor: o projeto é lualatex (fontspec/polyglossia).  Se o editor chamar pdflatex, cai
% num preâmbulo equivalente com inputenc/babel — serve para conferir, não para o PDF final.
\documentclass[11pt,a5paper]{book}
\ifdefined\directlua
  \usepackage{fontspec}
  \usepackage{polyglossia}
  \setmainlanguage{portuguese}
\else\ifdefined\XeTeXrevision
  \usepackage{fontspec}
  \usepackage{polyglossia}
  \setmainlanguage{portuguese}
\else
  \usepackage[utf8]{inputenc}
  \usepackage[T1]{fontenc}
  \usepackage{lmodern}
  \usepackage[portuguese]{babel}
  \DeclareUnicodeCharacter{2192}{\ensuremath{\rightarrow}}
\fi\fi
% raiz do repo: compilado da raiz ou de volN/<parte>/ (dois níveis abaixo)
\IfFileExists{tex/preambulo.tex}{% Preâmbulo comum ao livro e à pré-visualização de fragmentos (tex/_preview.tex).

% Mancha em A5.  Sem isto vale o padrão da classe book (~105 mm de largura e margem
% externa de ~25 mm em frente e verso), largo demais para o formato.  Aqui: ~118 mm de
% largura, margem interna um pouco maior que a externa (encadernação).
\usepackage[a5paper,inner=16mm,outer=14mm,top=14mm,bottom=16mm,
            includehead,headsep=5mm,footskip=8mm]{geometry}

% Títulos de secção sem justificar (evita linhas estouradas em \Large bold, p. ex.
% «Demonstrabilidade») e um pouco de folga para enunciados de tese em negrito.
\usepackage{sectsty}
\allsectionsfont{\raggedright}
\setlength{\emergencystretch}{1.5em}

% Marca da página impressa do original: \pg{286} imprime [286] no fluxo do texto.
\newcommand{\pg}[1]{\,{\footnotesize\textbf{[#1]}}\,}

% Virada de página do original DENTRO de uma nota (\pgn{422}); discreta, para não quebrar a nota.
\newcommand{\pgn}[1]{{\scriptsize\textbf{[#1]}}}

% Asterisco do original (divisão de assunto, sem título): \asterisco centrado em linha própria,
% também válido no meio do parágrafo ou dentro de nota (\footnote{... \asterisco ...}).
\newcommand{\asterisco}{\par\medskip\centerline{\footnotesize$\ast$}\medskip\noindent}

% Gravura recortada do fac-símile (<figure> do TEI; tools/tei2latex.py): fora de float, no ponto
% exato do original.  \gravura[detalhes]{arquivo.png}{fração da largura do texto}{legenda}:
% legenda principal centrada ("FIGURA VIII. — Côrte…"); os detalhes opcionais (<p> do <figure>:
% "A, hemispherio…; — B, …") vão abaixo, em corpo menor e justificados.  Ambos podem ser vazios.
\usepackage{graphicx}
\newcommand{\gravura}[4][]{\par\medskip\begingroup\centering
  \includegraphics[width=#3\linewidth,height=0.75\textheight,keepaspectratio]{#2}\par
  \if\relax\detokenize{#4}\relax\else\smallskip{\small #4\par}\fi
  \if\relax\detokenize{#1}\relax\else\smallskip{\footnotesize\leftskip=1em\rightskip=1em
    \parfillskip=0pt plus 1fil\noindent #1\par}\fi\endgroup\medskip\noindent}

% Citações latinas de S. Tomás (o OCR erra muito aqui: sempre conferir na imagem).
\newenvironment{latim}{\begin{quote}\itshape}{\end{quote}}

% --- Estrutura do autor (promovido de "Artigo III.tex") -----------------------
% enunciado de tese numerada que corre com o texto (segue-se travessão)
\newcommand{\tese}[2]{\bigskip\noindent\textbf{#1.\quad #2}\;\textemdash\;}
% enunciado de tese numerada em linha própria (seguem-se as subdivisões a), b)...)
\newcommand{\teseB}[2]{\bigskip\noindent\textbf{#1.\quad #2}\par}
% subdivisão a), b), c)... com rubrica em itálico
\newcommand{\sub}[2]{\medskip\noindent#1)\quad\emph{#2},\;\textemdash\;}
% subdivisão a), b), c)... sem rubrica
\newcommand{\subp}[1]{\medskip\noindent#1)\quad}

% --- Notas: no pé (padrão) ou correndo no corpo do texto ---------------------
% Passar  \def\notascorridas{}  no -usepretex faz cada nota correr no corpo, em
% corpo menor, num bloco recuado que interrompe a leitura até acabar.  Ex.:
%   latexmk -lualatex -usepretex='\def\fragmento{vol2/.../01_cap_y.tex}\def\notascorridas{}' tex/_preview.tex
\newif\ifnotascorridas
\ifdefined\notascorridas \notascorridastrue \fi
\ifnotascorridas
  \renewcommand\footnote[1]{%
    \stepcounter{footnote}\textsuperscript{\thefootnote}% marca no ponto do texto
    \par
    {\footnotesize
     \leftskip=2em \rightskip=2em \parindent=1em
     \addvspace{\medskipamount}%
     \noindent\textit{Nota \thefootnote.}\enspace #1\par
     \addvspace{\medskipamount}}%
  }
\fi

% --- Sem notas: \def\semnotas{} no -usepretex suprime todas as notas ---------
% (nem a chamada no texto nem o conteúdo).  Usado por  ./compila.sh --sem-notas .
\ifdefined\semnotas
  \renewcommand\footnote[2][]{}%
  \renewcommand\footnotemark[1][]{}%
  \renewcommand\footnotetext[2][]{}%
\fi

% peças pré-textuais (nó pretextual; p@type em web/obra/esquema.py)
\newcommand{\preTituloObra}[1]{\par\vspace{1em}{\centering\Large\bfseries\scshape #1\par}\vspace{1em}}
\newcommand{\preCentro}[1]{\par{\centering #1\par}}
\newcommand{\preEnderecamento}[1]{\par\vspace{1em}{\centering\small\scshape #1\par}}
\newcommand{\preSaudacao}[1]{\par\medskip\noindent #1\par}
\newcommand{\preDataLocal}[1]{\par\medskip\noindent\emph{#1}\par}
\newcommand{\preFecho}[1]{\par\smallskip{\raggedleft\scshape #1\par}}
}{% Preâmbulo comum ao livro e à pré-visualização de fragmentos (tex/_preview.tex).

% Mancha em A5.  Sem isto vale o padrão da classe book (~105 mm de largura e margem
% externa de ~25 mm em frente e verso), largo demais para o formato.  Aqui: ~118 mm de
% largura, margem interna um pouco maior que a externa (encadernação).
\usepackage[a5paper,inner=16mm,outer=14mm,top=14mm,bottom=16mm,
            includehead,headsep=5mm,footskip=8mm]{geometry}

% Títulos de secção sem justificar (evita linhas estouradas em \Large bold, p. ex.
% «Demonstrabilidade») e um pouco de folga para enunciados de tese em negrito.
\usepackage{sectsty}
\allsectionsfont{\raggedright}
\setlength{\emergencystretch}{1.5em}

% Marca da página impressa do original: \pg{286} imprime [286] no fluxo do texto.
\newcommand{\pg}[1]{\,{\footnotesize\textbf{[#1]}}\,}

% Virada de página do original DENTRO de uma nota (\pgn{422}); discreta, para não quebrar a nota.
\newcommand{\pgn}[1]{{\scriptsize\textbf{[#1]}}}

% Asterisco do original (divisão de assunto, sem título): \asterisco centrado em linha própria,
% também válido no meio do parágrafo ou dentro de nota (\footnote{... \asterisco ...}).
\newcommand{\asterisco}{\par\medskip\centerline{\footnotesize$\ast$}\medskip\noindent}

% Gravura recortada do fac-símile (<figure> do TEI; tools/tei2latex.py): fora de float, no ponto
% exato do original.  \gravura[detalhes]{arquivo.png}{fração da largura do texto}{legenda}:
% legenda principal centrada ("FIGURA VIII. — Côrte…"); os detalhes opcionais (<p> do <figure>:
% "A, hemispherio…; — B, …") vão abaixo, em corpo menor e justificados.  Ambos podem ser vazios.
\usepackage{graphicx}
\newcommand{\gravura}[4][]{\par\medskip\begingroup\centering
  \includegraphics[width=#3\linewidth,height=0.75\textheight,keepaspectratio]{#2}\par
  \if\relax\detokenize{#4}\relax\else\smallskip{\small #4\par}\fi
  \if\relax\detokenize{#1}\relax\else\smallskip{\footnotesize\leftskip=1em\rightskip=1em
    \parfillskip=0pt plus 1fil\noindent #1\par}\fi\endgroup\medskip\noindent}

% Citações latinas de S. Tomás (o OCR erra muito aqui: sempre conferir na imagem).
\newenvironment{latim}{\begin{quote}\itshape}{\end{quote}}

% --- Estrutura do autor (promovido de "Artigo III.tex") -----------------------
% enunciado de tese numerada que corre com o texto (segue-se travessão)
\newcommand{\tese}[2]{\bigskip\noindent\textbf{#1.\quad #2}\;\textemdash\;}
% enunciado de tese numerada em linha própria (seguem-se as subdivisões a), b)...)
\newcommand{\teseB}[2]{\bigskip\noindent\textbf{#1.\quad #2}\par}
% subdivisão a), b), c)... com rubrica em itálico
\newcommand{\sub}[2]{\medskip\noindent#1)\quad\emph{#2},\;\textemdash\;}
% subdivisão a), b), c)... sem rubrica
\newcommand{\subp}[1]{\medskip\noindent#1)\quad}

% --- Notas: no pé (padrão) ou correndo no corpo do texto ---------------------
% Passar  \def\notascorridas{}  no -usepretex faz cada nota correr no corpo, em
% corpo menor, num bloco recuado que interrompe a leitura até acabar.  Ex.:
%   latexmk -lualatex -usepretex='\def\fragmento{vol2/.../01_cap_y.tex}\def\notascorridas{}' tex/_preview.tex
\newif\ifnotascorridas
\ifdefined\notascorridas \notascorridastrue \fi
\ifnotascorridas
  \renewcommand\footnote[1]{%
    \stepcounter{footnote}\textsuperscript{\thefootnote}% marca no ponto do texto
    \par
    {\footnotesize
     \leftskip=2em \rightskip=2em \parindent=1em
     \addvspace{\medskipamount}%
     \noindent\textit{Nota \thefootnote.}\enspace #1\par
     \addvspace{\medskipamount}}%
  }
\fi

% --- Sem notas: \def\semnotas{} no -usepretex suprime todas as notas ---------
% (nem a chamada no texto nem o conteúdo).  Usado por  ./compila.sh --sem-notas .
\ifdefined\semnotas
  \renewcommand\footnote[2][]{}%
  \renewcommand\footnotemark[1][]{}%
  \renewcommand\footnotetext[2][]{}%
\fi

% peças pré-textuais (nó pretextual; p@type em web/obra/esquema.py)
\newcommand{\preTituloObra}[1]{\par\vspace{1em}{\centering\Large\bfseries\scshape #1\par}\vspace{1em}}
\newcommand{\preCentro}[1]{\par{\centering #1\par}}
\newcommand{\preEnderecamento}[1]{\par\vspace{1em}{\centering\small\scshape #1\par}}
\newcommand{\preSaudacao}[1]{\par\medskip\noindent #1\par}
\newcommand{\preDataLocal}[1]{\par\medskip\noindent\emph{#1}\par}
\newcommand{\preFecho}[1]{\par\smallskip{\raggedleft\scshape #1\par}}
}
\def\fimavulso{\end{document}}
\begin{document}
}\fi
%
% e termina com  \csname fimavulso\endcsname .  Dentro do livro ou do _preview, \pg já existe
% e nada disto roda; \fimavulso fica indefinido e o \csname vira \relax.
%
% Motor: o projeto é lualatex (fontspec/polyglossia).  Se o editor chamar pdflatex, cai
% num preâmbulo equivalente com inputenc/babel — serve para conferir, não para o PDF final.
\documentclass[11pt,a5paper]{book}
\ifdefined\directlua
  \usepackage{fontspec}
  \usepackage{polyglossia}
  \setmainlanguage{portuguese}
\else\ifdefined\XeTeXrevision
  \usepackage{fontspec}
  \usepackage{polyglossia}
  \setmainlanguage{portuguese}
\else
  \usepackage[utf8]{inputenc}
  \usepackage[T1]{fontenc}
  \usepackage{lmodern}
  \usepackage[portuguese]{babel}
  \DeclareUnicodeCharacter{2192}{\ensuremath{\rightarrow}}
\fi\fi
% raiz do repo: compilado da raiz ou de volN/<parte>/ (dois níveis abaixo)
\IfFileExists{tex/preambulo.tex}{% Preâmbulo comum ao livro e à pré-visualização de fragmentos (tex/_preview.tex).

% Mancha em A5.  Sem isto vale o padrão da classe book (~105 mm de largura e margem
% externa de ~25 mm em frente e verso), largo demais para o formato.  Aqui: ~118 mm de
% largura, margem interna um pouco maior que a externa (encadernação).
\usepackage[a5paper,inner=16mm,outer=14mm,top=14mm,bottom=16mm,
            includehead,headsep=5mm,footskip=8mm]{geometry}

% Títulos de secção sem justificar (evita linhas estouradas em \Large bold, p. ex.
% «Demonstrabilidade») e um pouco de folga para enunciados de tese em negrito.
\usepackage{sectsty}
\allsectionsfont{\raggedright}
\setlength{\emergencystretch}{1.5em}

% Marca da página impressa do original: \pg{286} imprime [286] no fluxo do texto.
\newcommand{\pg}[1]{\,{\footnotesize\textbf{[#1]}}\,}

% Virada de página do original DENTRO de uma nota (\pgn{422}); discreta, para não quebrar a nota.
\newcommand{\pgn}[1]{{\scriptsize\textbf{[#1]}}}

% Asterisco do original (divisão de assunto, sem título): \asterisco centrado em linha própria,
% também válido no meio do parágrafo ou dentro de nota (\footnote{... \asterisco ...}).
\newcommand{\asterisco}{\par\medskip\centerline{\footnotesize$\ast$}\medskip\noindent}

% Gravura recortada do fac-símile (<figure> do TEI; tools/tei2latex.py): fora de float, no ponto
% exato do original.  \gravura[detalhes]{arquivo.png}{fração da largura do texto}{legenda}:
% legenda principal centrada ("FIGURA VIII. — Côrte…"); os detalhes opcionais (<p> do <figure>:
% "A, hemispherio…; — B, …") vão abaixo, em corpo menor e justificados.  Ambos podem ser vazios.
\usepackage{graphicx}
\newcommand{\gravura}[4][]{\par\medskip\begingroup\centering
  \includegraphics[width=#3\linewidth,height=0.75\textheight,keepaspectratio]{#2}\par
  \if\relax\detokenize{#4}\relax\else\smallskip{\small #4\par}\fi
  \if\relax\detokenize{#1}\relax\else\smallskip{\footnotesize\leftskip=1em\rightskip=1em
    \parfillskip=0pt plus 1fil\noindent #1\par}\fi\endgroup\medskip\noindent}

% Citações latinas de S. Tomás (o OCR erra muito aqui: sempre conferir na imagem).
\newenvironment{latim}{\begin{quote}\itshape}{\end{quote}}

% --- Estrutura do autor (promovido de "Artigo III.tex") -----------------------
% enunciado de tese numerada que corre com o texto (segue-se travessão)
\newcommand{\tese}[2]{\bigskip\noindent\textbf{#1.\quad #2}\;\textemdash\;}
% enunciado de tese numerada em linha própria (seguem-se as subdivisões a), b)...)
\newcommand{\teseB}[2]{\bigskip\noindent\textbf{#1.\quad #2}\par}
% subdivisão a), b), c)... com rubrica em itálico
\newcommand{\sub}[2]{\medskip\noindent#1)\quad\emph{#2},\;\textemdash\;}
% subdivisão a), b), c)... sem rubrica
\newcommand{\subp}[1]{\medskip\noindent#1)\quad}

% --- Notas: no pé (padrão) ou correndo no corpo do texto ---------------------
% Passar  \def\notascorridas{}  no -usepretex faz cada nota correr no corpo, em
% corpo menor, num bloco recuado que interrompe a leitura até acabar.  Ex.:
%   latexmk -lualatex -usepretex='\def\fragmento{vol2/.../01_cap_y.tex}\def\notascorridas{}' tex/_preview.tex
\newif\ifnotascorridas
\ifdefined\notascorridas \notascorridastrue \fi
\ifnotascorridas
  \renewcommand\footnote[1]{%
    \stepcounter{footnote}\textsuperscript{\thefootnote}% marca no ponto do texto
    \par
    {\footnotesize
     \leftskip=2em \rightskip=2em \parindent=1em
     \addvspace{\medskipamount}%
     \noindent\textit{Nota \thefootnote.}\enspace #1\par
     \addvspace{\medskipamount}}%
  }
\fi

% --- Sem notas: \def\semnotas{} no -usepretex suprime todas as notas ---------
% (nem a chamada no texto nem o conteúdo).  Usado por  ./compila.sh --sem-notas .
\ifdefined\semnotas
  \renewcommand\footnote[2][]{}%
  \renewcommand\footnotemark[1][]{}%
  \renewcommand\footnotetext[2][]{}%
\fi

% peças pré-textuais (nó pretextual; p@type em web/obra/esquema.py)
\newcommand{\preTituloObra}[1]{\par\vspace{1em}{\centering\Large\bfseries\scshape #1\par}\vspace{1em}}
\newcommand{\preCentro}[1]{\par{\centering #1\par}}
\newcommand{\preEnderecamento}[1]{\par\vspace{1em}{\centering\small\scshape #1\par}}
\newcommand{\preSaudacao}[1]{\par\medskip\noindent #1\par}
\newcommand{\preDataLocal}[1]{\par\medskip\noindent\emph{#1}\par}
\newcommand{\preFecho}[1]{\par\smallskip{\raggedleft\scshape #1\par}}
}{% Preâmbulo comum ao livro e à pré-visualização de fragmentos (tex/_preview.tex).

% Mancha em A5.  Sem isto vale o padrão da classe book (~105 mm de largura e margem
% externa de ~25 mm em frente e verso), largo demais para o formato.  Aqui: ~118 mm de
% largura, margem interna um pouco maior que a externa (encadernação).
\usepackage[a5paper,inner=16mm,outer=14mm,top=14mm,bottom=16mm,
            includehead,headsep=5mm,footskip=8mm]{geometry}

% Títulos de secção sem justificar (evita linhas estouradas em \Large bold, p. ex.
% «Demonstrabilidade») e um pouco de folga para enunciados de tese em negrito.
\usepackage{sectsty}
\allsectionsfont{\raggedright}
\setlength{\emergencystretch}{1.5em}

% Marca da página impressa do original: \pg{286} imprime [286] no fluxo do texto.
\newcommand{\pg}[1]{\,{\footnotesize\textbf{[#1]}}\,}

% Virada de página do original DENTRO de uma nota (\pgn{422}); discreta, para não quebrar a nota.
\newcommand{\pgn}[1]{{\scriptsize\textbf{[#1]}}}

% Asterisco do original (divisão de assunto, sem título): \asterisco centrado em linha própria,
% também válido no meio do parágrafo ou dentro de nota (\footnote{... \asterisco ...}).
\newcommand{\asterisco}{\par\medskip\centerline{\footnotesize$\ast$}\medskip\noindent}

% Gravura recortada do fac-símile (<figure> do TEI; tools/tei2latex.py): fora de float, no ponto
% exato do original.  \gravura[detalhes]{arquivo.png}{fração da largura do texto}{legenda}:
% legenda principal centrada ("FIGURA VIII. — Côrte…"); os detalhes opcionais (<p> do <figure>:
% "A, hemispherio…; — B, …") vão abaixo, em corpo menor e justificados.  Ambos podem ser vazios.
\usepackage{graphicx}
\newcommand{\gravura}[4][]{\par\medskip\begingroup\centering
  \includegraphics[width=#3\linewidth,height=0.75\textheight,keepaspectratio]{#2}\par
  \if\relax\detokenize{#4}\relax\else\smallskip{\small #4\par}\fi
  \if\relax\detokenize{#1}\relax\else\smallskip{\footnotesize\leftskip=1em\rightskip=1em
    \parfillskip=0pt plus 1fil\noindent #1\par}\fi\endgroup\medskip\noindent}

% Citações latinas de S. Tomás (o OCR erra muito aqui: sempre conferir na imagem).
\newenvironment{latim}{\begin{quote}\itshape}{\end{quote}}

% --- Estrutura do autor (promovido de "Artigo III.tex") -----------------------
% enunciado de tese numerada que corre com o texto (segue-se travessão)
\newcommand{\tese}[2]{\bigskip\noindent\textbf{#1.\quad #2}\;\textemdash\;}
% enunciado de tese numerada em linha própria (seguem-se as subdivisões a), b)...)
\newcommand{\teseB}[2]{\bigskip\noindent\textbf{#1.\quad #2}\par}
% subdivisão a), b), c)... com rubrica em itálico
\newcommand{\sub}[2]{\medskip\noindent#1)\quad\emph{#2},\;\textemdash\;}
% subdivisão a), b), c)... sem rubrica
\newcommand{\subp}[1]{\medskip\noindent#1)\quad}

% --- Notas: no pé (padrão) ou correndo no corpo do texto ---------------------
% Passar  \def\notascorridas{}  no -usepretex faz cada nota correr no corpo, em
% corpo menor, num bloco recuado que interrompe a leitura até acabar.  Ex.:
%   latexmk -lualatex -usepretex='\def\fragmento{vol2/.../01_cap_y.tex}\def\notascorridas{}' tex/_preview.tex
\newif\ifnotascorridas
\ifdefined\notascorridas \notascorridastrue \fi
\ifnotascorridas
  \renewcommand\footnote[1]{%
    \stepcounter{footnote}\textsuperscript{\thefootnote}% marca no ponto do texto
    \par
    {\footnotesize
     \leftskip=2em \rightskip=2em \parindent=1em
     \addvspace{\medskipamount}%
     \noindent\textit{Nota \thefootnote.}\enspace #1\par
     \addvspace{\medskipamount}}%
  }
\fi

% --- Sem notas: \def\semnotas{} no -usepretex suprime todas as notas ---------
% (nem a chamada no texto nem o conteúdo).  Usado por  ./compila.sh --sem-notas .
\ifdefined\semnotas
  \renewcommand\footnote[2][]{}%
  \renewcommand\footnotemark[1][]{}%
  \renewcommand\footnotetext[2][]{}%
\fi

% peças pré-textuais (nó pretextual; p@type em web/obra/esquema.py)
\newcommand{\preTituloObra}[1]{\par\vspace{1em}{\centering\Large\bfseries\scshape #1\par}\vspace{1em}}
\newcommand{\preCentro}[1]{\par{\centering #1\par}}
\newcommand{\preEnderecamento}[1]{\par\vspace{1em}{\centering\small\scshape #1\par}}
\newcommand{\preSaudacao}[1]{\par\medskip\noindent #1\par}
\newcommand{\preDataLocal}[1]{\par\medskip\noindent\emph{#1}\par}
\newcommand{\preFecho}[1]{\par\smallskip{\raggedleft\scshape #1\par}}
}
\def\fimavulso{\end{document}}
\begin{document}
}{% Compilação AVULSA de um fragmento (volN/.../<stem>.tex) — p. ex. pelo botão do editor,
% de dentro da pasta do nó.  O fragmento começa com
%
%   \ifdefined\pg\else\IfFileExists{tex/avulso.tex}{% Compilação AVULSA de um fragmento (volN/.../<stem>.tex) — p. ex. pelo botão do editor,
% de dentro da pasta do nó.  O fragmento começa com
%
%   \ifdefined\pg\else\IfFileExists{tex/avulso.tex}{% Compilação AVULSA de um fragmento (volN/.../<stem>.tex) — p. ex. pelo botão do editor,
% de dentro da pasta do nó.  O fragmento começa com
%
%   \ifdefined\pg\else\IfFileExists{tex/avulso.tex}{\input{tex/avulso}}{\input{../../tex/avulso}}\fi
%
% e termina com  \csname fimavulso\endcsname .  Dentro do livro ou do _preview, \pg já existe
% e nada disto roda; \fimavulso fica indefinido e o \csname vira \relax.
%
% Motor: o projeto é lualatex (fontspec/polyglossia).  Se o editor chamar pdflatex, cai
% num preâmbulo equivalente com inputenc/babel — serve para conferir, não para o PDF final.
\documentclass[11pt,a5paper]{book}
\ifdefined\directlua
  \usepackage{fontspec}
  \usepackage{polyglossia}
  \setmainlanguage{portuguese}
\else\ifdefined\XeTeXrevision
  \usepackage{fontspec}
  \usepackage{polyglossia}
  \setmainlanguage{portuguese}
\else
  \usepackage[utf8]{inputenc}
  \usepackage[T1]{fontenc}
  \usepackage{lmodern}
  \usepackage[portuguese]{babel}
  \DeclareUnicodeCharacter{2192}{\ensuremath{\rightarrow}}
\fi\fi
% raiz do repo: compilado da raiz ou de volN/<parte>/ (dois níveis abaixo)
\IfFileExists{tex/preambulo.tex}{\input{tex/preambulo}}{\input{../../tex/preambulo}}
\def\fimavulso{\end{document}}
\begin{document}
}{% Compilação AVULSA de um fragmento (volN/.../<stem>.tex) — p. ex. pelo botão do editor,
% de dentro da pasta do nó.  O fragmento começa com
%
%   \ifdefined\pg\else\IfFileExists{tex/avulso.tex}{\input{tex/avulso}}{\input{../../tex/avulso}}\fi
%
% e termina com  \csname fimavulso\endcsname .  Dentro do livro ou do _preview, \pg já existe
% e nada disto roda; \fimavulso fica indefinido e o \csname vira \relax.
%
% Motor: o projeto é lualatex (fontspec/polyglossia).  Se o editor chamar pdflatex, cai
% num preâmbulo equivalente com inputenc/babel — serve para conferir, não para o PDF final.
\documentclass[11pt,a5paper]{book}
\ifdefined\directlua
  \usepackage{fontspec}
  \usepackage{polyglossia}
  \setmainlanguage{portuguese}
\else\ifdefined\XeTeXrevision
  \usepackage{fontspec}
  \usepackage{polyglossia}
  \setmainlanguage{portuguese}
\else
  \usepackage[utf8]{inputenc}
  \usepackage[T1]{fontenc}
  \usepackage{lmodern}
  \usepackage[portuguese]{babel}
  \DeclareUnicodeCharacter{2192}{\ensuremath{\rightarrow}}
\fi\fi
% raiz do repo: compilado da raiz ou de volN/<parte>/ (dois níveis abaixo)
\IfFileExists{tex/preambulo.tex}{\input{tex/preambulo}}{\input{../../tex/preambulo}}
\def\fimavulso{\end{document}}
\begin{document}
}\fi
%
% e termina com  \csname fimavulso\endcsname .  Dentro do livro ou do _preview, \pg já existe
% e nada disto roda; \fimavulso fica indefinido e o \csname vira \relax.
%
% Motor: o projeto é lualatex (fontspec/polyglossia).  Se o editor chamar pdflatex, cai
% num preâmbulo equivalente com inputenc/babel — serve para conferir, não para o PDF final.
\documentclass[11pt,a5paper]{book}
\ifdefined\directlua
  \usepackage{fontspec}
  \usepackage{polyglossia}
  \setmainlanguage{portuguese}
\else\ifdefined\XeTeXrevision
  \usepackage{fontspec}
  \usepackage{polyglossia}
  \setmainlanguage{portuguese}
\else
  \usepackage[utf8]{inputenc}
  \usepackage[T1]{fontenc}
  \usepackage{lmodern}
  \usepackage[portuguese]{babel}
  \DeclareUnicodeCharacter{2192}{\ensuremath{\rightarrow}}
\fi\fi
% raiz do repo: compilado da raiz ou de volN/<parte>/ (dois níveis abaixo)
\IfFileExists{tex/preambulo.tex}{% Preâmbulo comum ao livro e à pré-visualização de fragmentos (tex/_preview.tex).

% Mancha em A5.  Sem isto vale o padrão da classe book (~105 mm de largura e margem
% externa de ~25 mm em frente e verso), largo demais para o formato.  Aqui: ~118 mm de
% largura, margem interna um pouco maior que a externa (encadernação).
\usepackage[a5paper,inner=16mm,outer=14mm,top=14mm,bottom=16mm,
            includehead,headsep=5mm,footskip=8mm]{geometry}

% Títulos de secção sem justificar (evita linhas estouradas em \Large bold, p. ex.
% «Demonstrabilidade») e um pouco de folga para enunciados de tese em negrito.
\usepackage{sectsty}
\allsectionsfont{\raggedright}
\setlength{\emergencystretch}{1.5em}

% Marca da página impressa do original: \pg{286} imprime [286] no fluxo do texto.
\newcommand{\pg}[1]{\,{\footnotesize\textbf{[#1]}}\,}

% Virada de página do original DENTRO de uma nota (\pgn{422}); discreta, para não quebrar a nota.
\newcommand{\pgn}[1]{{\scriptsize\textbf{[#1]}}}

% Asterisco do original (divisão de assunto, sem título): \asterisco centrado em linha própria,
% também válido no meio do parágrafo ou dentro de nota (\footnote{... \asterisco ...}).
\newcommand{\asterisco}{\par\medskip\centerline{\footnotesize$\ast$}\medskip\noindent}

% Gravura recortada do fac-símile (<figure> do TEI; tools/tei2latex.py): fora de float, no ponto
% exato do original.  \gravura[detalhes]{arquivo.png}{fração da largura do texto}{legenda}:
% legenda principal centrada ("FIGURA VIII. — Côrte…"); os detalhes opcionais (<p> do <figure>:
% "A, hemispherio…; — B, …") vão abaixo, em corpo menor e justificados.  Ambos podem ser vazios.
\usepackage{graphicx}
\newcommand{\gravura}[4][]{\par\medskip\begingroup\centering
  \includegraphics[width=#3\linewidth,height=0.75\textheight,keepaspectratio]{#2}\par
  \if\relax\detokenize{#4}\relax\else\smallskip{\small #4\par}\fi
  \if\relax\detokenize{#1}\relax\else\smallskip{\footnotesize\leftskip=1em\rightskip=1em
    \parfillskip=0pt plus 1fil\noindent #1\par}\fi\endgroup\medskip\noindent}

% Citações latinas de S. Tomás (o OCR erra muito aqui: sempre conferir na imagem).
\newenvironment{latim}{\begin{quote}\itshape}{\end{quote}}

% --- Estrutura do autor (promovido de "Artigo III.tex") -----------------------
% enunciado de tese numerada que corre com o texto (segue-se travessão)
\newcommand{\tese}[2]{\bigskip\noindent\textbf{#1.\quad #2}\;\textemdash\;}
% enunciado de tese numerada em linha própria (seguem-se as subdivisões a), b)...)
\newcommand{\teseB}[2]{\bigskip\noindent\textbf{#1.\quad #2}\par}
% subdivisão a), b), c)... com rubrica em itálico
\newcommand{\sub}[2]{\medskip\noindent#1)\quad\emph{#2},\;\textemdash\;}
% subdivisão a), b), c)... sem rubrica
\newcommand{\subp}[1]{\medskip\noindent#1)\quad}

% --- Notas: no pé (padrão) ou correndo no corpo do texto ---------------------
% Passar  \def\notascorridas{}  no -usepretex faz cada nota correr no corpo, em
% corpo menor, num bloco recuado que interrompe a leitura até acabar.  Ex.:
%   latexmk -lualatex -usepretex='\def\fragmento{vol2/.../01_cap_y.tex}\def\notascorridas{}' tex/_preview.tex
\newif\ifnotascorridas
\ifdefined\notascorridas \notascorridastrue \fi
\ifnotascorridas
  \renewcommand\footnote[1]{%
    \stepcounter{footnote}\textsuperscript{\thefootnote}% marca no ponto do texto
    \par
    {\footnotesize
     \leftskip=2em \rightskip=2em \parindent=1em
     \addvspace{\medskipamount}%
     \noindent\textit{Nota \thefootnote.}\enspace #1\par
     \addvspace{\medskipamount}}%
  }
\fi

% --- Sem notas: \def\semnotas{} no -usepretex suprime todas as notas ---------
% (nem a chamada no texto nem o conteúdo).  Usado por  ./compila.sh --sem-notas .
\ifdefined\semnotas
  \renewcommand\footnote[2][]{}%
  \renewcommand\footnotemark[1][]{}%
  \renewcommand\footnotetext[2][]{}%
\fi

% peças pré-textuais (nó pretextual; p@type em web/obra/esquema.py)
\newcommand{\preTituloObra}[1]{\par\vspace{1em}{\centering\Large\bfseries\scshape #1\par}\vspace{1em}}
\newcommand{\preCentro}[1]{\par{\centering #1\par}}
\newcommand{\preEnderecamento}[1]{\par\vspace{1em}{\centering\small\scshape #1\par}}
\newcommand{\preSaudacao}[1]{\par\medskip\noindent #1\par}
\newcommand{\preDataLocal}[1]{\par\medskip\noindent\emph{#1}\par}
\newcommand{\preFecho}[1]{\par\smallskip{\raggedleft\scshape #1\par}}
}{% Preâmbulo comum ao livro e à pré-visualização de fragmentos (tex/_preview.tex).

% Mancha em A5.  Sem isto vale o padrão da classe book (~105 mm de largura e margem
% externa de ~25 mm em frente e verso), largo demais para o formato.  Aqui: ~118 mm de
% largura, margem interna um pouco maior que a externa (encadernação).
\usepackage[a5paper,inner=16mm,outer=14mm,top=14mm,bottom=16mm,
            includehead,headsep=5mm,footskip=8mm]{geometry}

% Títulos de secção sem justificar (evita linhas estouradas em \Large bold, p. ex.
% «Demonstrabilidade») e um pouco de folga para enunciados de tese em negrito.
\usepackage{sectsty}
\allsectionsfont{\raggedright}
\setlength{\emergencystretch}{1.5em}

% Marca da página impressa do original: \pg{286} imprime [286] no fluxo do texto.
\newcommand{\pg}[1]{\,{\footnotesize\textbf{[#1]}}\,}

% Virada de página do original DENTRO de uma nota (\pgn{422}); discreta, para não quebrar a nota.
\newcommand{\pgn}[1]{{\scriptsize\textbf{[#1]}}}

% Asterisco do original (divisão de assunto, sem título): \asterisco centrado em linha própria,
% também válido no meio do parágrafo ou dentro de nota (\footnote{... \asterisco ...}).
\newcommand{\asterisco}{\par\medskip\centerline{\footnotesize$\ast$}\medskip\noindent}

% Gravura recortada do fac-símile (<figure> do TEI; tools/tei2latex.py): fora de float, no ponto
% exato do original.  \gravura[detalhes]{arquivo.png}{fração da largura do texto}{legenda}:
% legenda principal centrada ("FIGURA VIII. — Côrte…"); os detalhes opcionais (<p> do <figure>:
% "A, hemispherio…; — B, …") vão abaixo, em corpo menor e justificados.  Ambos podem ser vazios.
\usepackage{graphicx}
\newcommand{\gravura}[4][]{\par\medskip\begingroup\centering
  \includegraphics[width=#3\linewidth,height=0.75\textheight,keepaspectratio]{#2}\par
  \if\relax\detokenize{#4}\relax\else\smallskip{\small #4\par}\fi
  \if\relax\detokenize{#1}\relax\else\smallskip{\footnotesize\leftskip=1em\rightskip=1em
    \parfillskip=0pt plus 1fil\noindent #1\par}\fi\endgroup\medskip\noindent}

% Citações latinas de S. Tomás (o OCR erra muito aqui: sempre conferir na imagem).
\newenvironment{latim}{\begin{quote}\itshape}{\end{quote}}

% --- Estrutura do autor (promovido de "Artigo III.tex") -----------------------
% enunciado de tese numerada que corre com o texto (segue-se travessão)
\newcommand{\tese}[2]{\bigskip\noindent\textbf{#1.\quad #2}\;\textemdash\;}
% enunciado de tese numerada em linha própria (seguem-se as subdivisões a), b)...)
\newcommand{\teseB}[2]{\bigskip\noindent\textbf{#1.\quad #2}\par}
% subdivisão a), b), c)... com rubrica em itálico
\newcommand{\sub}[2]{\medskip\noindent#1)\quad\emph{#2},\;\textemdash\;}
% subdivisão a), b), c)... sem rubrica
\newcommand{\subp}[1]{\medskip\noindent#1)\quad}

% --- Notas: no pé (padrão) ou correndo no corpo do texto ---------------------
% Passar  \def\notascorridas{}  no -usepretex faz cada nota correr no corpo, em
% corpo menor, num bloco recuado que interrompe a leitura até acabar.  Ex.:
%   latexmk -lualatex -usepretex='\def\fragmento{vol2/.../01_cap_y.tex}\def\notascorridas{}' tex/_preview.tex
\newif\ifnotascorridas
\ifdefined\notascorridas \notascorridastrue \fi
\ifnotascorridas
  \renewcommand\footnote[1]{%
    \stepcounter{footnote}\textsuperscript{\thefootnote}% marca no ponto do texto
    \par
    {\footnotesize
     \leftskip=2em \rightskip=2em \parindent=1em
     \addvspace{\medskipamount}%
     \noindent\textit{Nota \thefootnote.}\enspace #1\par
     \addvspace{\medskipamount}}%
  }
\fi

% --- Sem notas: \def\semnotas{} no -usepretex suprime todas as notas ---------
% (nem a chamada no texto nem o conteúdo).  Usado por  ./compila.sh --sem-notas .
\ifdefined\semnotas
  \renewcommand\footnote[2][]{}%
  \renewcommand\footnotemark[1][]{}%
  \renewcommand\footnotetext[2][]{}%
\fi

% peças pré-textuais (nó pretextual; p@type em web/obra/esquema.py)
\newcommand{\preTituloObra}[1]{\par\vspace{1em}{\centering\Large\bfseries\scshape #1\par}\vspace{1em}}
\newcommand{\preCentro}[1]{\par{\centering #1\par}}
\newcommand{\preEnderecamento}[1]{\par\vspace{1em}{\centering\small\scshape #1\par}}
\newcommand{\preSaudacao}[1]{\par\medskip\noindent #1\par}
\newcommand{\preDataLocal}[1]{\par\medskip\noindent\emph{#1}\par}
\newcommand{\preFecho}[1]{\par\smallskip{\raggedleft\scshape #1\par}}
}
\def\fimavulso{\end{document}}
\begin{document}
}{% Compilação AVULSA de um fragmento (volN/.../<stem>.tex) — p. ex. pelo botão do editor,
% de dentro da pasta do nó.  O fragmento começa com
%
%   \ifdefined\pg\else\IfFileExists{tex/avulso.tex}{% Compilação AVULSA de um fragmento (volN/.../<stem>.tex) — p. ex. pelo botão do editor,
% de dentro da pasta do nó.  O fragmento começa com
%
%   \ifdefined\pg\else\IfFileExists{tex/avulso.tex}{\input{tex/avulso}}{\input{../../tex/avulso}}\fi
%
% e termina com  \csname fimavulso\endcsname .  Dentro do livro ou do _preview, \pg já existe
% e nada disto roda; \fimavulso fica indefinido e o \csname vira \relax.
%
% Motor: o projeto é lualatex (fontspec/polyglossia).  Se o editor chamar pdflatex, cai
% num preâmbulo equivalente com inputenc/babel — serve para conferir, não para o PDF final.
\documentclass[11pt,a5paper]{book}
\ifdefined\directlua
  \usepackage{fontspec}
  \usepackage{polyglossia}
  \setmainlanguage{portuguese}
\else\ifdefined\XeTeXrevision
  \usepackage{fontspec}
  \usepackage{polyglossia}
  \setmainlanguage{portuguese}
\else
  \usepackage[utf8]{inputenc}
  \usepackage[T1]{fontenc}
  \usepackage{lmodern}
  \usepackage[portuguese]{babel}
  \DeclareUnicodeCharacter{2192}{\ensuremath{\rightarrow}}
\fi\fi
% raiz do repo: compilado da raiz ou de volN/<parte>/ (dois níveis abaixo)
\IfFileExists{tex/preambulo.tex}{\input{tex/preambulo}}{\input{../../tex/preambulo}}
\def\fimavulso{\end{document}}
\begin{document}
}{% Compilação AVULSA de um fragmento (volN/.../<stem>.tex) — p. ex. pelo botão do editor,
% de dentro da pasta do nó.  O fragmento começa com
%
%   \ifdefined\pg\else\IfFileExists{tex/avulso.tex}{\input{tex/avulso}}{\input{../../tex/avulso}}\fi
%
% e termina com  \csname fimavulso\endcsname .  Dentro do livro ou do _preview, \pg já existe
% e nada disto roda; \fimavulso fica indefinido e o \csname vira \relax.
%
% Motor: o projeto é lualatex (fontspec/polyglossia).  Se o editor chamar pdflatex, cai
% num preâmbulo equivalente com inputenc/babel — serve para conferir, não para o PDF final.
\documentclass[11pt,a5paper]{book}
\ifdefined\directlua
  \usepackage{fontspec}
  \usepackage{polyglossia}
  \setmainlanguage{portuguese}
\else\ifdefined\XeTeXrevision
  \usepackage{fontspec}
  \usepackage{polyglossia}
  \setmainlanguage{portuguese}
\else
  \usepackage[utf8]{inputenc}
  \usepackage[T1]{fontenc}
  \usepackage{lmodern}
  \usepackage[portuguese]{babel}
  \DeclareUnicodeCharacter{2192}{\ensuremath{\rightarrow}}
\fi\fi
% raiz do repo: compilado da raiz ou de volN/<parte>/ (dois níveis abaixo)
\IfFileExists{tex/preambulo.tex}{\input{tex/preambulo}}{\input{../../tex/preambulo}}
\def\fimavulso{\end{document}}
\begin{document}
}\fi
%
% e termina com  \csname fimavulso\endcsname .  Dentro do livro ou do _preview, \pg já existe
% e nada disto roda; \fimavulso fica indefinido e o \csname vira \relax.
%
% Motor: o projeto é lualatex (fontspec/polyglossia).  Se o editor chamar pdflatex, cai
% num preâmbulo equivalente com inputenc/babel — serve para conferir, não para o PDF final.
\documentclass[11pt,a5paper]{book}
\ifdefined\directlua
  \usepackage{fontspec}
  \usepackage{polyglossia}
  \setmainlanguage{portuguese}
\else\ifdefined\XeTeXrevision
  \usepackage{fontspec}
  \usepackage{polyglossia}
  \setmainlanguage{portuguese}
\else
  \usepackage[utf8]{inputenc}
  \usepackage[T1]{fontenc}
  \usepackage{lmodern}
  \usepackage[portuguese]{babel}
  \DeclareUnicodeCharacter{2192}{\ensuremath{\rightarrow}}
\fi\fi
% raiz do repo: compilado da raiz ou de volN/<parte>/ (dois níveis abaixo)
\IfFileExists{tex/preambulo.tex}{% Preâmbulo comum ao livro e à pré-visualização de fragmentos (tex/_preview.tex).

% Mancha em A5.  Sem isto vale o padrão da classe book (~105 mm de largura e margem
% externa de ~25 mm em frente e verso), largo demais para o formato.  Aqui: ~118 mm de
% largura, margem interna um pouco maior que a externa (encadernação).
\usepackage[a5paper,inner=16mm,outer=14mm,top=14mm,bottom=16mm,
            includehead,headsep=5mm,footskip=8mm]{geometry}

% Títulos de secção sem justificar (evita linhas estouradas em \Large bold, p. ex.
% «Demonstrabilidade») e um pouco de folga para enunciados de tese em negrito.
\usepackage{sectsty}
\allsectionsfont{\raggedright}
\setlength{\emergencystretch}{1.5em}

% Marca da página impressa do original: \pg{286} imprime [286] no fluxo do texto.
\newcommand{\pg}[1]{\,{\footnotesize\textbf{[#1]}}\,}

% Virada de página do original DENTRO de uma nota (\pgn{422}); discreta, para não quebrar a nota.
\newcommand{\pgn}[1]{{\scriptsize\textbf{[#1]}}}

% Asterisco do original (divisão de assunto, sem título): \asterisco centrado em linha própria,
% também válido no meio do parágrafo ou dentro de nota (\footnote{... \asterisco ...}).
\newcommand{\asterisco}{\par\medskip\centerline{\footnotesize$\ast$}\medskip\noindent}

% Gravura recortada do fac-símile (<figure> do TEI; tools/tei2latex.py): fora de float, no ponto
% exato do original.  \gravura[detalhes]{arquivo.png}{fração da largura do texto}{legenda}:
% legenda principal centrada ("FIGURA VIII. — Côrte…"); os detalhes opcionais (<p> do <figure>:
% "A, hemispherio…; — B, …") vão abaixo, em corpo menor e justificados.  Ambos podem ser vazios.
\usepackage{graphicx}
\newcommand{\gravura}[4][]{\par\medskip\begingroup\centering
  \includegraphics[width=#3\linewidth,height=0.75\textheight,keepaspectratio]{#2}\par
  \if\relax\detokenize{#4}\relax\else\smallskip{\small #4\par}\fi
  \if\relax\detokenize{#1}\relax\else\smallskip{\footnotesize\leftskip=1em\rightskip=1em
    \parfillskip=0pt plus 1fil\noindent #1\par}\fi\endgroup\medskip\noindent}

% Citações latinas de S. Tomás (o OCR erra muito aqui: sempre conferir na imagem).
\newenvironment{latim}{\begin{quote}\itshape}{\end{quote}}

% --- Estrutura do autor (promovido de "Artigo III.tex") -----------------------
% enunciado de tese numerada que corre com o texto (segue-se travessão)
\newcommand{\tese}[2]{\bigskip\noindent\textbf{#1.\quad #2}\;\textemdash\;}
% enunciado de tese numerada em linha própria (seguem-se as subdivisões a), b)...)
\newcommand{\teseB}[2]{\bigskip\noindent\textbf{#1.\quad #2}\par}
% subdivisão a), b), c)... com rubrica em itálico
\newcommand{\sub}[2]{\medskip\noindent#1)\quad\emph{#2},\;\textemdash\;}
% subdivisão a), b), c)... sem rubrica
\newcommand{\subp}[1]{\medskip\noindent#1)\quad}

% --- Notas: no pé (padrão) ou correndo no corpo do texto ---------------------
% Passar  \def\notascorridas{}  no -usepretex faz cada nota correr no corpo, em
% corpo menor, num bloco recuado que interrompe a leitura até acabar.  Ex.:
%   latexmk -lualatex -usepretex='\def\fragmento{vol2/.../01_cap_y.tex}\def\notascorridas{}' tex/_preview.tex
\newif\ifnotascorridas
\ifdefined\notascorridas \notascorridastrue \fi
\ifnotascorridas
  \renewcommand\footnote[1]{%
    \stepcounter{footnote}\textsuperscript{\thefootnote}% marca no ponto do texto
    \par
    {\footnotesize
     \leftskip=2em \rightskip=2em \parindent=1em
     \addvspace{\medskipamount}%
     \noindent\textit{Nota \thefootnote.}\enspace #1\par
     \addvspace{\medskipamount}}%
  }
\fi

% --- Sem notas: \def\semnotas{} no -usepretex suprime todas as notas ---------
% (nem a chamada no texto nem o conteúdo).  Usado por  ./compila.sh --sem-notas .
\ifdefined\semnotas
  \renewcommand\footnote[2][]{}%
  \renewcommand\footnotemark[1][]{}%
  \renewcommand\footnotetext[2][]{}%
\fi

% peças pré-textuais (nó pretextual; p@type em web/obra/esquema.py)
\newcommand{\preTituloObra}[1]{\par\vspace{1em}{\centering\Large\bfseries\scshape #1\par}\vspace{1em}}
\newcommand{\preCentro}[1]{\par{\centering #1\par}}
\newcommand{\preEnderecamento}[1]{\par\vspace{1em}{\centering\small\scshape #1\par}}
\newcommand{\preSaudacao}[1]{\par\medskip\noindent #1\par}
\newcommand{\preDataLocal}[1]{\par\medskip\noindent\emph{#1}\par}
\newcommand{\preFecho}[1]{\par\smallskip{\raggedleft\scshape #1\par}}
}{% Preâmbulo comum ao livro e à pré-visualização de fragmentos (tex/_preview.tex).

% Mancha em A5.  Sem isto vale o padrão da classe book (~105 mm de largura e margem
% externa de ~25 mm em frente e verso), largo demais para o formato.  Aqui: ~118 mm de
% largura, margem interna um pouco maior que a externa (encadernação).
\usepackage[a5paper,inner=16mm,outer=14mm,top=14mm,bottom=16mm,
            includehead,headsep=5mm,footskip=8mm]{geometry}

% Títulos de secção sem justificar (evita linhas estouradas em \Large bold, p. ex.
% «Demonstrabilidade») e um pouco de folga para enunciados de tese em negrito.
\usepackage{sectsty}
\allsectionsfont{\raggedright}
\setlength{\emergencystretch}{1.5em}

% Marca da página impressa do original: \pg{286} imprime [286] no fluxo do texto.
\newcommand{\pg}[1]{\,{\footnotesize\textbf{[#1]}}\,}

% Virada de página do original DENTRO de uma nota (\pgn{422}); discreta, para não quebrar a nota.
\newcommand{\pgn}[1]{{\scriptsize\textbf{[#1]}}}

% Asterisco do original (divisão de assunto, sem título): \asterisco centrado em linha própria,
% também válido no meio do parágrafo ou dentro de nota (\footnote{... \asterisco ...}).
\newcommand{\asterisco}{\par\medskip\centerline{\footnotesize$\ast$}\medskip\noindent}

% Gravura recortada do fac-símile (<figure> do TEI; tools/tei2latex.py): fora de float, no ponto
% exato do original.  \gravura[detalhes]{arquivo.png}{fração da largura do texto}{legenda}:
% legenda principal centrada ("FIGURA VIII. — Côrte…"); os detalhes opcionais (<p> do <figure>:
% "A, hemispherio…; — B, …") vão abaixo, em corpo menor e justificados.  Ambos podem ser vazios.
\usepackage{graphicx}
\newcommand{\gravura}[4][]{\par\medskip\begingroup\centering
  \includegraphics[width=#3\linewidth,height=0.75\textheight,keepaspectratio]{#2}\par
  \if\relax\detokenize{#4}\relax\else\smallskip{\small #4\par}\fi
  \if\relax\detokenize{#1}\relax\else\smallskip{\footnotesize\leftskip=1em\rightskip=1em
    \parfillskip=0pt plus 1fil\noindent #1\par}\fi\endgroup\medskip\noindent}

% Citações latinas de S. Tomás (o OCR erra muito aqui: sempre conferir na imagem).
\newenvironment{latim}{\begin{quote}\itshape}{\end{quote}}

% --- Estrutura do autor (promovido de "Artigo III.tex") -----------------------
% enunciado de tese numerada que corre com o texto (segue-se travessão)
\newcommand{\tese}[2]{\bigskip\noindent\textbf{#1.\quad #2}\;\textemdash\;}
% enunciado de tese numerada em linha própria (seguem-se as subdivisões a), b)...)
\newcommand{\teseB}[2]{\bigskip\noindent\textbf{#1.\quad #2}\par}
% subdivisão a), b), c)... com rubrica em itálico
\newcommand{\sub}[2]{\medskip\noindent#1)\quad\emph{#2},\;\textemdash\;}
% subdivisão a), b), c)... sem rubrica
\newcommand{\subp}[1]{\medskip\noindent#1)\quad}

% --- Notas: no pé (padrão) ou correndo no corpo do texto ---------------------
% Passar  \def\notascorridas{}  no -usepretex faz cada nota correr no corpo, em
% corpo menor, num bloco recuado que interrompe a leitura até acabar.  Ex.:
%   latexmk -lualatex -usepretex='\def\fragmento{vol2/.../01_cap_y.tex}\def\notascorridas{}' tex/_preview.tex
\newif\ifnotascorridas
\ifdefined\notascorridas \notascorridastrue \fi
\ifnotascorridas
  \renewcommand\footnote[1]{%
    \stepcounter{footnote}\textsuperscript{\thefootnote}% marca no ponto do texto
    \par
    {\footnotesize
     \leftskip=2em \rightskip=2em \parindent=1em
     \addvspace{\medskipamount}%
     \noindent\textit{Nota \thefootnote.}\enspace #1\par
     \addvspace{\medskipamount}}%
  }
\fi

% --- Sem notas: \def\semnotas{} no -usepretex suprime todas as notas ---------
% (nem a chamada no texto nem o conteúdo).  Usado por  ./compila.sh --sem-notas .
\ifdefined\semnotas
  \renewcommand\footnote[2][]{}%
  \renewcommand\footnotemark[1][]{}%
  \renewcommand\footnotetext[2][]{}%
\fi

% peças pré-textuais (nó pretextual; p@type em web/obra/esquema.py)
\newcommand{\preTituloObra}[1]{\par\vspace{1em}{\centering\Large\bfseries\scshape #1\par}\vspace{1em}}
\newcommand{\preCentro}[1]{\par{\centering #1\par}}
\newcommand{\preEnderecamento}[1]{\par\vspace{1em}{\centering\small\scshape #1\par}}
\newcommand{\preSaudacao}[1]{\par\medskip\noindent #1\par}
\newcommand{\preDataLocal}[1]{\par\medskip\noindent\emph{#1}\par}
\newcommand{\preFecho}[1]{\par\smallskip{\raggedleft\scshape #1\par}}
}
\def\fimavulso{\end{document}}
\begin{document}
}\fi
%
% e termina com  \csname fimavulso\endcsname .  Dentro do livro ou do _preview, \pg já existe
% e nada disto roda; \fimavulso fica indefinido e o \csname vira \relax.
%
% Motor: o projeto é lualatex (fontspec/polyglossia).  Se o editor chamar pdflatex, cai
% num preâmbulo equivalente com inputenc/babel — serve para conferir, não para o PDF final.
\documentclass[11pt,a5paper]{book}
\ifdefined\directlua
  \usepackage{fontspec}
  \usepackage{polyglossia}
  \setmainlanguage{portuguese}
\else\ifdefined\XeTeXrevision
  \usepackage{fontspec}
  \usepackage{polyglossia}
  \setmainlanguage{portuguese}
\else
  \usepackage[utf8]{inputenc}
  \usepackage[T1]{fontenc}
  \usepackage{lmodern}
  \usepackage[portuguese]{babel}
  \DeclareUnicodeCharacter{2192}{\ensuremath{\rightarrow}}
\fi\fi
% raiz do repo: compilado da raiz ou de volN/<parte>/ (dois níveis abaixo)
\IfFileExists{tex/preambulo.tex}{% Preâmbulo comum ao livro e à pré-visualização de fragmentos (tex/_preview.tex).

% Mancha em A5.  Sem isto vale o padrão da classe book (~105 mm de largura e margem
% externa de ~25 mm em frente e verso), largo demais para o formato.  Aqui: ~118 mm de
% largura, margem interna um pouco maior que a externa (encadernação).
\usepackage[a5paper,inner=16mm,outer=14mm,top=14mm,bottom=16mm,
            includehead,headsep=5mm,footskip=8mm]{geometry}

% Títulos de secção sem justificar (evita linhas estouradas em \Large bold, p. ex.
% «Demonstrabilidade») e um pouco de folga para enunciados de tese em negrito.
\usepackage{sectsty}
\allsectionsfont{\raggedright}
\setlength{\emergencystretch}{1.5em}

% Marca da página impressa do original: \pg{286} imprime [286] no fluxo do texto.
\newcommand{\pg}[1]{\,{\footnotesize\textbf{[#1]}}\,}

% Virada de página do original DENTRO de uma nota (\pgn{422}); discreta, para não quebrar a nota.
\newcommand{\pgn}[1]{{\scriptsize\textbf{[#1]}}}

% Asterisco do original (divisão de assunto, sem título): \asterisco centrado em linha própria,
% também válido no meio do parágrafo ou dentro de nota (\footnote{... \asterisco ...}).
\newcommand{\asterisco}{\par\medskip\centerline{\footnotesize$\ast$}\medskip\noindent}

% Gravura recortada do fac-símile (<figure> do TEI; tools/tei2latex.py): fora de float, no ponto
% exato do original.  \gravura[detalhes]{arquivo.png}{fração da largura do texto}{legenda}:
% legenda principal centrada ("FIGURA VIII. — Côrte…"); os detalhes opcionais (<p> do <figure>:
% "A, hemispherio…; — B, …") vão abaixo, em corpo menor e justificados.  Ambos podem ser vazios.
\usepackage{graphicx}
\newcommand{\gravura}[4][]{\par\medskip\begingroup\centering
  \includegraphics[width=#3\linewidth,height=0.75\textheight,keepaspectratio]{#2}\par
  \if\relax\detokenize{#4}\relax\else\smallskip{\small #4\par}\fi
  \if\relax\detokenize{#1}\relax\else\smallskip{\footnotesize\leftskip=1em\rightskip=1em
    \parfillskip=0pt plus 1fil\noindent #1\par}\fi\endgroup\medskip\noindent}

% Citações latinas de S. Tomás (o OCR erra muito aqui: sempre conferir na imagem).
\newenvironment{latim}{\begin{quote}\itshape}{\end{quote}}

% --- Estrutura do autor (promovido de "Artigo III.tex") -----------------------
% enunciado de tese numerada que corre com o texto (segue-se travessão)
\newcommand{\tese}[2]{\bigskip\noindent\textbf{#1.\quad #2}\;\textemdash\;}
% enunciado de tese numerada em linha própria (seguem-se as subdivisões a), b)...)
\newcommand{\teseB}[2]{\bigskip\noindent\textbf{#1.\quad #2}\par}
% subdivisão a), b), c)... com rubrica em itálico
\newcommand{\sub}[2]{\medskip\noindent#1)\quad\emph{#2},\;\textemdash\;}
% subdivisão a), b), c)... sem rubrica
\newcommand{\subp}[1]{\medskip\noindent#1)\quad}

% --- Notas: no pé (padrão) ou correndo no corpo do texto ---------------------
% Passar  \def\notascorridas{}  no -usepretex faz cada nota correr no corpo, em
% corpo menor, num bloco recuado que interrompe a leitura até acabar.  Ex.:
%   latexmk -lualatex -usepretex='\def\fragmento{vol2/.../01_cap_y.tex}\def\notascorridas{}' tex/_preview.tex
\newif\ifnotascorridas
\ifdefined\notascorridas \notascorridastrue \fi
\ifnotascorridas
  \renewcommand\footnote[1]{%
    \stepcounter{footnote}\textsuperscript{\thefootnote}% marca no ponto do texto
    \par
    {\footnotesize
     \leftskip=2em \rightskip=2em \parindent=1em
     \addvspace{\medskipamount}%
     \noindent\textit{Nota \thefootnote.}\enspace #1\par
     \addvspace{\medskipamount}}%
  }
\fi

% --- Sem notas: \def\semnotas{} no -usepretex suprime todas as notas ---------
% (nem a chamada no texto nem o conteúdo).  Usado por  ./compila.sh --sem-notas .
\ifdefined\semnotas
  \renewcommand\footnote[2][]{}%
  \renewcommand\footnotemark[1][]{}%
  \renewcommand\footnotetext[2][]{}%
\fi

% peças pré-textuais (nó pretextual; p@type em web/obra/esquema.py)
\newcommand{\preTituloObra}[1]{\par\vspace{1em}{\centering\Large\bfseries\scshape #1\par}\vspace{1em}}
\newcommand{\preCentro}[1]{\par{\centering #1\par}}
\newcommand{\preEnderecamento}[1]{\par\vspace{1em}{\centering\small\scshape #1\par}}
\newcommand{\preSaudacao}[1]{\par\medskip\noindent #1\par}
\newcommand{\preDataLocal}[1]{\par\medskip\noindent\emph{#1}\par}
\newcommand{\preFecho}[1]{\par\smallskip{\raggedleft\scshape #1\par}}
}{% Preâmbulo comum ao livro e à pré-visualização de fragmentos (tex/_preview.tex).

% Mancha em A5.  Sem isto vale o padrão da classe book (~105 mm de largura e margem
% externa de ~25 mm em frente e verso), largo demais para o formato.  Aqui: ~118 mm de
% largura, margem interna um pouco maior que a externa (encadernação).
\usepackage[a5paper,inner=16mm,outer=14mm,top=14mm,bottom=16mm,
            includehead,headsep=5mm,footskip=8mm]{geometry}

% Títulos de secção sem justificar (evita linhas estouradas em \Large bold, p. ex.
% «Demonstrabilidade») e um pouco de folga para enunciados de tese em negrito.
\usepackage{sectsty}
\allsectionsfont{\raggedright}
\setlength{\emergencystretch}{1.5em}

% Marca da página impressa do original: \pg{286} imprime [286] no fluxo do texto.
\newcommand{\pg}[1]{\,{\footnotesize\textbf{[#1]}}\,}

% Virada de página do original DENTRO de uma nota (\pgn{422}); discreta, para não quebrar a nota.
\newcommand{\pgn}[1]{{\scriptsize\textbf{[#1]}}}

% Asterisco do original (divisão de assunto, sem título): \asterisco centrado em linha própria,
% também válido no meio do parágrafo ou dentro de nota (\footnote{... \asterisco ...}).
\newcommand{\asterisco}{\par\medskip\centerline{\footnotesize$\ast$}\medskip\noindent}

% Gravura recortada do fac-símile (<figure> do TEI; tools/tei2latex.py): fora de float, no ponto
% exato do original.  \gravura[detalhes]{arquivo.png}{fração da largura do texto}{legenda}:
% legenda principal centrada ("FIGURA VIII. — Côrte…"); os detalhes opcionais (<p> do <figure>:
% "A, hemispherio…; — B, …") vão abaixo, em corpo menor e justificados.  Ambos podem ser vazios.
\usepackage{graphicx}
\newcommand{\gravura}[4][]{\par\medskip\begingroup\centering
  \includegraphics[width=#3\linewidth,height=0.75\textheight,keepaspectratio]{#2}\par
  \if\relax\detokenize{#4}\relax\else\smallskip{\small #4\par}\fi
  \if\relax\detokenize{#1}\relax\else\smallskip{\footnotesize\leftskip=1em\rightskip=1em
    \parfillskip=0pt plus 1fil\noindent #1\par}\fi\endgroup\medskip\noindent}

% Citações latinas de S. Tomás (o OCR erra muito aqui: sempre conferir na imagem).
\newenvironment{latim}{\begin{quote}\itshape}{\end{quote}}

% --- Estrutura do autor (promovido de "Artigo III.tex") -----------------------
% enunciado de tese numerada que corre com o texto (segue-se travessão)
\newcommand{\tese}[2]{\bigskip\noindent\textbf{#1.\quad #2}\;\textemdash\;}
% enunciado de tese numerada em linha própria (seguem-se as subdivisões a), b)...)
\newcommand{\teseB}[2]{\bigskip\noindent\textbf{#1.\quad #2}\par}
% subdivisão a), b), c)... com rubrica em itálico
\newcommand{\sub}[2]{\medskip\noindent#1)\quad\emph{#2},\;\textemdash\;}
% subdivisão a), b), c)... sem rubrica
\newcommand{\subp}[1]{\medskip\noindent#1)\quad}

% --- Notas: no pé (padrão) ou correndo no corpo do texto ---------------------
% Passar  \def\notascorridas{}  no -usepretex faz cada nota correr no corpo, em
% corpo menor, num bloco recuado que interrompe a leitura até acabar.  Ex.:
%   latexmk -lualatex -usepretex='\def\fragmento{vol2/.../01_cap_y.tex}\def\notascorridas{}' tex/_preview.tex
\newif\ifnotascorridas
\ifdefined\notascorridas \notascorridastrue \fi
\ifnotascorridas
  \renewcommand\footnote[1]{%
    \stepcounter{footnote}\textsuperscript{\thefootnote}% marca no ponto do texto
    \par
    {\footnotesize
     \leftskip=2em \rightskip=2em \parindent=1em
     \addvspace{\medskipamount}%
     \noindent\textit{Nota \thefootnote.}\enspace #1\par
     \addvspace{\medskipamount}}%
  }
\fi

% --- Sem notas: \def\semnotas{} no -usepretex suprime todas as notas ---------
% (nem a chamada no texto nem o conteúdo).  Usado por  ./compila.sh --sem-notas .
\ifdefined\semnotas
  \renewcommand\footnote[2][]{}%
  \renewcommand\footnotemark[1][]{}%
  \renewcommand\footnotetext[2][]{}%
\fi

% peças pré-textuais (nó pretextual; p@type em web/obra/esquema.py)
\newcommand{\preTituloObra}[1]{\par\vspace{1em}{\centering\Large\bfseries\scshape #1\par}\vspace{1em}}
\newcommand{\preCentro}[1]{\par{\centering #1\par}}
\newcommand{\preEnderecamento}[1]{\par\vspace{1em}{\centering\small\scshape #1\par}}
\newcommand{\preSaudacao}[1]{\par\medskip\noindent #1\par}
\newcommand{\preDataLocal}[1]{\par\medskip\noindent\emph{#1}\par}
\newcommand{\preFecho}[1]{\par\smallskip{\raggedleft\scshape #1\par}}
}
\def\fimavulso{\end{document}}
\begin{document}
}\fi
% PADRÃO-OURO vol. 1 (5.ª ed., 1927), bloco C — Metaphysica, Parte Segunda: Cosmologia,
% Prologo, Secção Primeira, Cap. I, Art. I.  Págs. impressas 577–580
% (vol1/paginas/0577.png … 0580.png).  Transcrito à mão da imagem; o OCR serviu só de
% rascunho.  Convenções em 062-065_logica_raciocinio.tex.  Aqui:
%   - a p. 577 é portada de parte (Metaphysica / Parte Segunda / Cosmologia / Prologo) e abre
%     sem resto de nota da página anterior;
%   - latim nas notas (Plinio; Rom. I, 19-20) e nomes próprios com « w » (Newton);
%   - a nota (1) da p. 580 acaba na p. 581;
%   - preservados: « causas supremus » (p. 577), « Liepig », « nas ontologica », « pelo
%     sentidos » (p. 581), « Creacianismo » (Summario).

\pg{577}

\part*{Metaphysica\\ Parte Segunda\\ Cosmologia}

\chapter*{Prologo}

\tese{1}{Cosmologia.} \emph{Cosmologia} é a \emph{sciencia, que trata das causas supremus
do mundo}. --- A \emph{Cosmologia} --- \emph{a}) é \emph{sciencia}; porque possue todos os
elementos, que constituem a sciencia; --- \emph{b}) trata das causas \emph{supremas},
porque a Philosophia estuda as mais altas causas dos entes; --- \emph{c}) trata das causas
supremas do \emph{mundo}, isto é, da universalidade das creaturas visiveis, contidas no
ceu e na terra, e conhecidas por meio dos sentidos\footnote{\emph{Cosmologia}, quanto ao
\emph{nome}, significa \emph{a sciencia do mundo}. \emph{Mundo} quer dizer
\emph{ornamento}. Disse \emph{Plinio}: «~Quem Graeci \emph{cosmon}, nomine
\emph{ornamenti}, appellavere, eum et nos a perfecta absolutaque elegantia \emph{mundum}
dicimus~» (\emph{Hist. nat.}, l. II, c. 4). --- A Cosmologia costuma chamar-se tambem
\emph{Physica racional}; porque trata das coisas sensiveis, não exclusivamente com o
auxilio da \emph{experiencia} (como a Physica \emph{experimental}), mas tambem com a luz
da \emph{razão}, baseada na mesma experiencia.}.

\tese{2}{Objecto da Cosmologia.} O \emph{objecto} da \emph{Cosmologia} é o \emph{mundo},
considerado nas suas \emph{causas supremas}. --- A \emph{Cosmologia}, pois, poderia
abranger não só os mineraes, os vegetaes e os animaes irracionaes, mas tambem o homem,
porque tambem este faz parte da universalidade das creaturas visiveis. --- Todavia, por
ser o homem composto tambem de uma \emph{substancia espiritual} e \emph{incorruptivel},
que o torna immensamente superior a todos os outros seres sensiveis, o tratado do
\emph{homem} forma uma parte distincta da Metaphysica, que é a
\emph{Anthropologia}\footnote{Os escriptores não concordam na determinação do
\emph{objecto} da Cosmologia. --- Alguns antigos e modernos (\emph{d'Alembert, Littré})
não distinguem a Cosmologia das sciencias physicas e mathematicas, pois nesta parte
estudam a \emph{uranologia}, a \emph{aerologia}, a \emph{geologia}, a
\emph{hydrologia}, etc. --- Outros identificam a Cosmologia com as outras partes da
Philosophia, estudando a \emph{essencia} dos corpos na Ontologia, o \emph{milagre} na
Logica, a \emph{origem} e os \emph{caracteres} do mundo na Theodicea. --- Outros tratam
na Cosmologia não só dos \emph{mineraes}, dos \emph{vegetaes} e dos \emph{animaes
irracionaes}, mas tambem do \emph{homem}. --- Outros finalmente dão ao \emph{objecto} da
Cosmologia a extensão que nós lhe attribuimos, reservando para o estudo acerca do
\emph{homem} um tratado especial.}.

\pg{578}
\tese{3}{A Cosmologia e as sciencias naturaes.} A \emph{Cosmologia} não só não impede,
mas até favorece o progresso das sciencias naturaes. --- Na verdade, ao passo que as
sciencias naturaes se limitam á observação das causas \emph{proximas} e
\emph{particulares} dos entes corporeos, a \emph{Cosmologia}, subindo mais alto,
investiga --- as leis \emph{universaes}, a que esses entes obedecem, --- os principios
\emph{supremos}, de que derivam ou de que são compostos, --- a \emph{ordem}, que
conservam entre si, e o \emph{fim}, para o qual tendem. D'este modo as sciencias
naturaes encontram na Cosmologia o criterio facil e seguro para as suas investigações, e
a razão ultima das suas descobertas\footnote{Por isso o estudo da \emph{Cosmologia}, ---
não só é \emph{util}, emquanto satisfaz plenamente ás legitimas aspirações da nossa
intelligencia, que deseja conhecer as causas supremas dos phenomenos sensiveis, --- mas é
tambem \emph{necessario} para nos elevarmos ao conhecimento de \textbf{Deus}, nosso
Creador e Senhor. Na verdade, \textbf{Deus}, invisivel por si mesmo aos nossos olhos,
quiz tornar-se visivel na ordem natural, por meio das coisas creadas, emquanto estas são
imagens e reflexos, ainda que immensamente pallidos, das adoraveis perfeições divinas.
Diz o Apostolo: «~\textbf{Deus} illis manifestavit. Invisibilia enim Ipsius a creatura
mundi, per ea quae facta sunt, intellecta, conspiciuntur; sempiterna quoque Ejus virtus et
divinitas~» (\emph{Rom.}, I, 19-20).

Entre a Philosophia e as sciencias naturaes tem existido até agora uma desconfiança
mutua, para não dizermos uma lucta aberta. Os cultores das sciencias naturaes desprezam
as verdades racionaes, julgando-as contrarias ás suas descobertas. Mas erram; não só
porque as verdades \emph{ontologicas}, que se encontram nas coisas, e as verdades
\emph{logicas}, que existem na intelligencia creada, são diversos raios do unico e
infinito sol de toda a verdade, que é \textbf{Deus}; mas tambem porque as verdades
\emph{logicas} se fundam nas \emph{ontologica}, e \emph{nunca} se provou que uma
descoberta \emph{certa} das sciencias naturaes fosse contraria a uma verdade racional.
Espiritos vastos e penetrantes, que alcançaram com o seu olhar o immenso campo da
philosophia e das sciencias, guiados não pelos prejuizos e pelas paixões mas pelo amor
sereno da verdade, confessaram que entre as verdades racionaes e as experimentaes existe
a mais perfeita harmonia. Citaremos, entre tantos, os nomes de \emph{Newton, Galileu,
Cuvier, Linneu, Kepler, Liepig, Flourens, Quatrefages, Moigno, Pasteur}. --- E
necessario, pois, harmonisar a sciencia racional e a sciencia experimental; porque, só
assim, teremos um \emph{completo} conhecimento das coisas. Quem só estuda a Philosophia
\emph{racional}, só conhece as causas \emph{supremas} das coisas, e este conhecimento é
\emph{imperfeito}; --- quem só estuda as sciencias experimentaes, só conhece as causas
\emph{proximas} das coisas, e este conhecimento é ainda mais imperfeito. Conhecer toda a
cadeia ou serie das causas, desde a mais proxima até á mais alta, é completa
sabedoria.}.

\tese{4}{Divisão da Cosmologia.} Dividimos a \emph{Cosmologia} em \emph{duas} secções. Na
primeira secção tratamos do \emph{mundo em geral}, da sua \emph{origem}, \emph{ordem} e
\emph{leis}, e da \emph{essencia} e das \emph{propriedades}\pg{579} dos \emph{corpos}; na
segunda, dos \emph{mineraes}, dos \emph{vegetaes}, e dos \emph{animaes irracionaes}. ---
A primeira secção constitue a \emph{Cosmologia geral}; a segunda secção, a
\emph{Cosmologia especial}.

\part*{Secção Primeira\\ Cosmologia geral}

\tese{5}{Divisão da Cosmologia geral.} Dividimos a Cosmologia \emph{geral} em tres
capitulos. No \emph{primeiro} tratamos da natureza e origem do mundo; no \emph{segundo},
da sua ordem e das suas leis; no \emph{terceiro}, da essencia e das propriedades dos
corpos.

\chapter*{Capitulo Primeiro\\[0.3em]\large Natureza e origem do mundo}

{\small\noindent\textsc{Summario}: --- Existencia real do mundo. --- Caracteres do mundo.
--- Origem do mundo: Materialismo. --- Origem do mundo: Pantheismo. --- Origem do mundo:
Creacianismo. --- Duração do mundo.\par}

\section*{Artigo I\\ Existencia real do mundo}

\tese{6}{Estado da questão.} A existencia real do mundo é uma verdade tão evidente, que,
para negal-a ou para duvidar d'ella, é necessario renunciar ao proprio senso commum. ---
Todavia alguns escriptores disseram que a idea, que temos do mundo, é mera ficção do
nosso espirito, porque o mundo não existe realmente. --- D'ahi a necessidade do presente
artigo\footnote{O systema, que nega a existencia \emph{real} do mundo, reduzindo-o á mera
idea, a uma esteril ficção do nosso espirito, chama-se \emph{idealismo}. --- O
\emph{idealismo} é \emph{vulgar} e \emph{transcendental}. O \emph{vulgar} nega a realidade
do mundo, mas admitte a realidade do sujeito pensante. O \emph{transcendental} nega toda a
realidade, e affirma que o mundo e o sujeito pensante são meras apparencias. --- O
idealismo \emph{vulgar} teve por defensores --- na antiguidade, os scepticos, --- e nos
tempos modernos, \emph{Jorge Berkeley, Schopenauer, Lange, Lotze}, e outros. --- O
\emph{transcendental} foi inventado por \emph{Kant}, e levado ao exaggero por
\emph{Fichte, Schelling} e \emph{Hegel}.

Deixando para mais tarde a refutação do idealismo \emph{transcendental}, neste artigo
limitamo-nos á critica do \emph{vulgar}.}.

\pg{580}
\teseB{7}{A idea, que temos do mundo, não é mera ficção da nossa intelligencia.}

\subp{a} Se a idea, que temos do mundo, fosse mera ficção da nossa intelligencia, o mundo
deveria apresentar-se ao nosso espirito como um ente \emph{simplicissimo} e
\emph{espiritual}; porque \emph{simplicissima} e \emph{espiritual} é a nossa
intelligencia, e toda a faculdade, que na sua operação não é determinada por um agente
extrinseco, produz um effeito \emph{semelhante} a si mesma. Ora o mundo apresenta-se ao
espirito como um ente \emph{composto} e \emph{material}. Logo a idea do mundo não é mera
ficção da nossa intelligencia.

\subp{b} Se a idea do \emph{mundo} fosse uma ficção da nossa intelligencia, o mundo
deveria representar-se-nos --- não só como existente no nosso espirito, porque é ahi que
está a idea, --- mas tambem de um modo \emph{abstracto}, porque abstracta é a idea. Ora o
mundo representa-se-nos como existente fóra do nosso espirito, e de um modo
\emph{concreto}. Logo a idea do mundo não é mera ficção da nossa intelligencia.

\subp{c} Se o mundo fosse mera ficção da nossa intelligencia, todas as mudanças, a que
elle está sujeito, deveriam ter a sua causa nas nossas operações intellectuaes. Ora a
experiencia attesta-nos que aquellas mudanças não dependem das nossas operações, mas
unica e exclusivamente das causas extrinsecas. Logo o mundo não é mera ficção da nossa
intelligencia\footnote{Entre tantos desvarios da intelligencia humana, o
\emph{idealismo} é um dos mais absurdos. Se o mundo é mera ficção da nossa intelligencia,
segue-se necessariamente --- que somos nós que formamos o modelo do mundo, --- que
existem tantos mundos, quantas são as intelligencias humanas, --- que a existencia do
mundo principia e acaba com a existencia de cada homem, --- que somos nós os verdadeiros
e unicos auctores d'esta ordem, d'esta variedade, d'esta belleza, que resplandece nas
creaturas, --- que somos nós que produzimos o sol, a lua, e as estrellas, e que damos a
vida ás plantas, a sensibilidade aos animaes, a intelligencia ao homem, --- que somos nós
os auctores de todas as pinturas, que contemplamos, de todos os livros, que lemos, de
todas as esculpturas, que admiramos, --- que somos nós que produzimos tudo o que nos
encanta, e o produzimos sem estudo, sem trabalho, e até sem consciencia! ---
Considerariam os idealistas estas consequencias?

Omittimos, porque facilmente se entendem, outras consequencias, e muito perniciosas, que
derivam do \emph{idealismo} para o campo da Fé e da Religião. Se os corpos não existem na
realidade, --- o Filho de \textbf{Deus} não pode ter assumido um corpo humano, e por isso
não pode ter nascido e morrido pela salvação do mundo, --- tudo o que a Fé ensina ácerca
da origem do mundo, da formação do homem, da instituição da Egreja e dos Sacramentos,
etc., tudo isto é uma illusão fundada sobre outra illusão.

O fundamento do idealismo \emph{vulgar} foi lançado por \emph{Descartes} e \emph{Locke},
quando disseram que \emph{só percebemos as nossas modificações e impressões subjectivas}.
--- Baseando-se neste principio, \emph{Condillac} concluiu que \emph{os corpos são o
conjuncto das sensações subjectivas}.

Os \emph{idealistas} cahiram neste erro, suppondo que podesse dar-se uma sensação
\emph{exclusivamente} subjectiva, isto é, uma sensação sem um objecto sentido; e assim
concluiram que só percebemos as nossas modificações subjectivas. Ora esta supposição e
esta conclusão são falsas. --- A supposição é falsa; porque a sensação, se é subjectiva
na \emph{essencia}, em quanto é uma operação immanente, é objectiva na \emph{origem} e na
\emph{representação}, em quanto é produzida por um objecto e representa um objecto. --- A
conclusão tambem é falsa; não só porque deriva de um principio falso, mas tambem porque
\emph{directa} e \emph{immediatamente} percebemos, pelo sentidos \emph{externos}, os
objectos \emph{reaes}, e não as modificações \emph{subjectivas}.}.

\csname fimavulso\endcsname
